% ================================
% THƯ MỤC CHÚ GIẢI
% ================================
\OpenGrid

\ParallelChapterStar
  {Annotated Bibliography}
  {Thư mục chú giải}
  {}

\TriPara
  {
    \EN{In this section, several citations from each chapter of the publication have been summarized to provide additional insights into, and support for, the ideas discussed.}
  }
  {
    \VI{Trong mục này, một số tài liệu được trích dẫn ở từng chương của ấn phẩm đã được tóm lược nhằm cung cấp thêm góc nhìn và cơ sở hỗ trợ cho những ý tưởng đã được thảo luận.}
  }
  {
    \NOTE{Không phải là danh mục tài liệu tham khảo đơn thuần, mà là tập hợp tài liệu được trích dẫn, đã được chắt lọc và chú giải nhằm làm sáng tỏ, củng cố luận điểm chính của toàn bộ ấn phẩm.

    Người đọc cần tầng hỗ trợ thứ hai để hiểu sâu hơn ý tưởng đã được trình bày trong từng chương, đặc biệt là ý tưởng liên quan đến luận và hiểu trong dạy học toán. Thư mục chú giải đóng vai trò cung cấp tầng hỗ trợ đó thông qua việc kết nối nội dung sách với nghiên cứu, quan điểm và công trình nền tảng.

    Điểm cần làm rõ là từ ``summarized'' ở đây không mang nghĩa rút gọn hình thức, mà là tóm lược tài liệu theo hướng làm nổi bật ý nghĩa sư phạm của chúng đối với chủ đề đang bàn. Điều này cho thấy trích dẫn được sử dụng có chủ đích, phục vụ lập luận, chứ không mang tính liệt kê.

    Về mặt cấu trúc sách, đoạn văn này đóng vai trò định vị: Thư mục chú giải là phần mở rộng học thuật, giúp người đọc, đặc biệt là giáo viên và người thiết kế chương trình, truy ngược về cơ sở lí luận của các quan điểm được nêu trong sách.}
  }

\ParallelSection
  {Chapter 1: Reasoning and Sense Making}
  {Chương 1: Luận và hiểu}
  {}

\TriPara
  {
    \EN{
      \begin{enumerate}
        \item American Diploma Project. \emph{Ready or Not: Creating a High School Diploma That Counts}. Washington, D.C.: Achieve, 2004.\vspace{.25\baselineskip}

          \par\noindent\hspace*{1em} This report strives to strengthen the connection between high school mathematics and English curricula and the skills and knowledge that high school graduates will need for either college or postsecondary work. Working closely with leaders of higher education as well as workforce leaders, the American Diploma Project describes in this report what students need to know on high school graduation so that they can be successful. Mathematical reasoning is highlighted as important for high school students because both in the workplace and in high school they will be called on to apply mathematical concepts from the classroom to new problems and situations.\vspace{1.35\baselineskip}

        \item Kilpatrick, Jeremy, Jane Swafford, and Bradford Findell, eds. \emph{Adding It Up: Helping Children Learn Mathematics}. Washington, D.C.: National Academy Press, 2001.\vspace{.25\baselineskip}

          \par\noindent\hspace*{1em} This National Research Council report synthesizes research and makes recommendations for the teaching and learning of mathematics for kindergarten through eighth-grade students. The report puts forth a model for describing mathematical proficiency that includes conceptual understanding, procedural fluency, strategic competence, adaptive reasoning, and productive disposition. The report describes how each of the five ``strands'' is a crucial component of students' success in school mathematics.\vspace{.35\baselineskip}

        \item Knuth, Eric J. ``Teachers' Conceptions of Proof in the Context of Secondary School Mathematics.'' \emph{Journal of Mathematics Teacher Education} 5, no. 1 (2002): 61--88.\vspace{.25\baselineskip}

          \par\noindent\hspace*{1em} In this study looking at secondary school mathematics teachers' conception of proof, Knuth presents a framework for understanding the role of proof in school mathematics. This framework draws on other research and provides a useful way of thinking about the purposes of proof in mathematics curricula. In terms of Knuth's framework, proof is a mechanism for verifying and explaining true statements, communicating mathematically, creating and discovering new mathematics, and organizing ideas so that they are part of a structure of axioms.\vspace{1.35\baselineskip}

        \item Yackel, Erna, and Gila Hanna. ``Reasoning and Proof.'' In \emph{A Research Companion to ``Principles and Standards for School Mathematics''}, edited by Jeremy Kilpatrick, W. Gary Martin, and Deborah Shifter, pp. 227--36. Reston, Va.: National Council of Teachers of Mathematics, 2003.\vspace{.25\baselineskip}

          \par\noindent\hspace*{1em} This chapter highlights research on reasoning and proof in K-12 education. The authors report research that supports the importance of reasoning in students' learning of mathematics. Further, they discuss the crucial role of proof in the field of mathematics and consequently in K-12 mathematics. Proof is more than just a means of verification of a statement; it also serves as a vehicle for explanation and systemization of observed relationships and patterns.
      \end{enumerate}}
  }
  {
    \VI{
      \begin{enumerate}
        \item American Diploma Project. \emph{Ready or Not: Creating a High School Diploma That Counts}. Washington, D.C.: Achieve, 2004.\vspace{.25\baselineskip}

          \par\noindent\hspace*{1em} Báo cáo này nỗ lực tăng cường kết nối giữa chương trình học toán và tiếng Anh ở bậc trung học phổ thông với kĩ năng và kiến thức mà học sinh tốt nghiệp trung học cần có để học đại học hoặc đi làm sau trung học. Thông qua hợp tác chặt chẽ với các nhà lãnh đạo giáo dục đại học cũng như lãnh đạo trong lực lượng lao động, American Diploma Project mô tả trong báo cáo này những điều học sinh cần biết khi tốt nghiệp trung học để có thể thành công. Luận toán học được nhấn mạnh là yếu tố quan trọng đối với học sinh trung học phổ thông, bởi cả trong môi trường làm việc lẫn học tập ở trường phổ thông, các em đều sẽ được yêu cầu vận dụng khái niệm toán học học trong lớp vào vấn đề và tình huống mới.\vspace{.35\baselineskip}

        \item Kilpatrick, Jeremy, Jane Swafford, and Bradford Findell, eds. \emph{Adding It Up: Helping Children Learn Mathematics}. Washington, D.C.: National Academy Press, 2001.\vspace{.25\baselineskip}

          \par\noindent\hspace*{1em} Báo cáo này của Hội đồng Nghiên cứu Quốc gia Hoa Kì tổng hợp nghiên cứu và đưa ra khuyến nghị cho việc dạy và học toán từ bậc mẫu giáo đến hết lớp tám. Báo cáo đề xuất mô hình mô tả năng lực toán học, bao gồm năm thành tố: hiểu biết khái niệm, thành thạo về thủ tục, năng lực chiến lược, luận thích ứng, và thái độ tích cực. Báo cáo cũng trình bày rằng mỗi trong năm ``thành tố'' này đều là bộ phận then chốt, góp phần quyết định sự thành công của học sinh trong việc học toán ở nhà trường.\vspace{.35\baselineskip}

        \item Knuth, Eric J. ``Teachers' Conceptions of Proof in the Context of Secondary School Mathematics.'' \emph{Journal of Mathematics Teacher Education} 5, no. 1 (2002): 61--88.\vspace{.25\baselineskip}

          \par\noindent\hspace*{1em} Trong nghiên cứu này về quan niệm của giáo viên toán bậc trung học đối với chứng minh, Knuth trình bày một khung lí thuyết nhằm giúp hiểu vai trò của chứng minh trong toán học nhà trường. Khung lí thuyết này dựa trên các nghiên cứu khác và cung cấp một cách hữu ích để suy nghĩ về mục đích của chứng minh trong chương trình học toán. Theo khung của Knuth, chứng minh là cơ chế để kiểm chứng và giải thích mệnh đề đúng, giao tiếp toán học, tạo ra và khám phá ý tưởng toán học mới, đồng thời tổ chức ý tưởng sao cho chúng trở thành một phần của cấu trúc tiên đề.\vspace{.35\baselineskip}

        \item Yackel, Erna, and Gila Hanna. ``Reasoning and Proof.'' In \emph{A Research Companion to ``Principles and Standards for School Mathematics''}, edited by Jeremy Kilpatrick, W. Gary Martin, and Deborah Shifter, pp. 227--36. Reston, Va.: National Council of Teachers of Mathematics, 2003.\vspace{.25\baselineskip}

          \par\noindent\hspace*{1em} Chương này nêu bật các nghiên cứu về luận và chứng minh trong giáo dục K--12\footnote{K--12 là cách gọi trong hệ thống giáo dục Hoa Kì, chỉ giáo dục từ mẫu giáo đến lớp 12, tức toàn bộ giai đoạn trước đại học. Trong bản dịch này, kí hiệu K--12 được giữ nguyên vì đây là tên gọi quy ước của hệ thống giáo dục Hoa Kì.}. Các tác giả trình bày những nghiên cứu ủng hộ tầm quan trọng của luận trong việc học toán của học sinh. Ngoài ra, họ thảo luận vai trò cốt yếu của chứng minh trong lĩnh vực toán học và do đó trong toán K--12. Chứng minh không chỉ là phương tiện kiểm chứng một mệnh đề; nó còn là phương tiện để giải thích và hệ thống hóa các quan hệ cùng mẫu hình quan sát được.
      \end{enumerate}}
  }{\NOTE{}}

\clearpage
\ParallelSection
  {Chapter 2: Reasoning Habits}
  {Chương 2: Thói quen luận}
  {}

\TriPara
  {
    \EN{
      \begin{enumerate}
        \item Cuoco, Al, E. Paul Goldenberg, and June Mark. ``Habits of Mind: An Organizing Principle for Mathematics Curriculum.'' \emph{Journal of Mathematical Behavior} 15 (December 1996): 375--402.\vspace{.25\baselineskip}

          \par\noindent\hspace*{1em} The authors make the case for emphasizing more than just mathematical content in high school mathematics. They propose that high school students need to develop mathematical habits of mind. These ways of thinking about mathematics, modeled after the work of mathematicians, will continue to be relevant to students even as content changes over time. The article describes in detail the habits of mind that high school students should be developing as part of their mathematics experience.\vspace{.35\baselineskip}

        \item Committee on the Undergraduate Program in Mathematics of the Mathematical Association of America (CUPM). \emph{Undergraduate Programs and Courses in the Mathematical Sciences: CUPM Curriculum Guide 2004}. Washington, D.C.: Mathematical Association of America, 2004.\vspace{.25\baselineskip}

          \par\noindent\hspace*{1em} The CUPM of the Mathematical Association of America is charged with making recommendations to colleges and universities about undergraduate mathematics education. Every ten years, on the basis of consultations with mathematicians and representatives from related disciplines, the committee updates and makes available the list of recommendations. Among other recommendations, the report points to the need for every undergraduate mathematics course to develop critical reasoning and mathematical habits of mind in students.\vspace*{.35\baselineskip}

        \item Garfield, Joan. ``The Challenge of Developing Statistical Reasoning.'' \emph{Journal of Statistics Education} 10, no. 3 (2002). http://www.amstat.org/publications/jse/v10n3/garfield.html.\vspace{.25\baselineskip}

          \par\noindent\hspace*{1em} The author defines and discusses research on statistical reasoning. She also describes a developmental model portraying five stages of statistical reasoning for students. In addition, she considers implications for assessment and teaching.\vspace{.35\baselineskip}

        \item Harel, Guershon, and Larry Sowder. ``Advanced Mathematical-Thinking at Any Age: Its Nature and Its Development.'' \emph{Mathematical Thinking and Learning} 7, no. 1 (2005): 27--50.\vspace{.25\baselineskip}

          \par\noindent\hspace*{1em} This article makes the case that advanced mathematical thinking is not just mathematical thinking about advanced mathematics but advanced thinking about mathematics that can occur as early as in elementary school. Students have many important opportunities to develop such ways of thinking in elementary and secondary school mathematics. The authors also describe hindrances, including particular instructional practices, to students developing advanced mathematical thinking.
      \end{enumerate}}
  }
  {
    \VI{
      \begin{enumerate}
        \item Cuoco, Al, E. Paul Goldenberg, and June Mark. ``Habits of Mind: An Organizing Principle for Mathematics Curriculum.'' \emph{Journal of Mathematical Behavior} 15 (December 1996): 375--402.\vspace{.25\baselineskip}

          \par\noindent\hspace*{1em} Các tác giả lập luận rằng toán trung học phổ thông không nên chỉ tập trung vào nội dung toán học. Họ đề xuất rằng học sinh trung học phổ thông cần phát triển thói quen tư duy toán học. Những cách suy nghĩ về toán này, được mô phỏng theo cách làm việc của nhà toán học, sẽ tiếp tục có ý nghĩa đối với học sinh ngay cả khi nội dung thay đổi theo thời gian. Bài báo mô tả chi tiết những thói quen tư duy mà học sinh trung học phổ thông nên phát triển như một phần trong trải nghiệm học toán.\vspace{.25\baselineskip}

        \item Committee on the Undergraduate Program in Mathematics of the Mathematical Association of America (CUPM). \emph{Undergraduate Programs and Courses in the Mathematical Sciences: CUPM Curriculum Guide 2004}. Washington, D.C.: Mathematical Association of America, 2004.\vspace{.25\baselineskip}

          \par\noindent\hspace*{1em} Ủy ban Chương trình Toán bậc đại học (CUPM) của Hiệp hội Toán học Hoa Kì được giao nhiệm vụ đưa ra khuyến nghị cho các trường cao đẳng và đại học về giáo dục toán bậc đại học. Cứ mười năm một lần, dựa trên tham vấn với các nhà toán học và đại diện từ các lĩnh vực liên quan, ủy ban cập nhật và công bố danh sách khuyến nghị. Trong số các khuyến nghị khác, báo cáo chỉ ra nhu cầu rằng mọi học phần toán ở bậc đại học đều phải phát triển luận phản biện và thói quen tư duy toán học cho sinh viên.\vspace*{.35\baselineskip}

        \item Garfield, Joan. ``The Challenge of Developing Statistical Reasoning.'' \emph{Journal of Statistics Education} 10, no. 3 (2002). http://www.amstat.org/publications/jse/v10n3/garfield.html.\vspace{.25\baselineskip}

          \par\noindent\hspace*{1em} Tác giả định nghĩa luận thống kê và thảo luận các nghiên cứu về chủ đề này. Bà cũng mô tả một mô hình phát triển gồm năm giai đoạn luận thống kê của học sinh. Ngoài ra, bà xem xét hàm ý đối với đánh giá và dạy học.\vspace{1.35\baselineskip}

        \item Harel, Guershon, and Larry Sowder. ``Advanced Mathematical-Thinking at Any Age: Its Nature and Its Development.'' \emph{Mathematical Thinking and Learning} 7, no. 1 (2005): 27--50.\vspace{.25\baselineskip}

          \par\noindent\hspace*{1em} Bài báo này lập luận rằng tư duy toán học nâng cao không chỉ là tư duy toán học về toán học nâng cao, mà là cách tư duy nâng cao về toán học, có thể xuất hiện ngay từ bậc tiểu học. Học sinh có nhiều cơ hội quan trọng để phát triển những cách tư duy như vậy trong toán ở bậc tiểu học và trung học. Các tác giả cũng mô tả những trở ngại, bao gồm một số thực hành dạy học cụ thể, đối với quá trình học sinh phát triển tư duy toán học nâng cao.
      \end{enumerate}}
  }{\NOTE{}}

\ParallelSection
  {Chapter 4: Reasoning with Number and Measurement}
  {Chương 4: Luận với số và đo lường}
  {}

\TriPara
  {
    \EN{
      \begin{enumerate}
        \item Carraher, David W., and Analucia D. Schliemann, ``Early Algebra and Algebraic Reasoning.'' In \emph{Second Handbook of Research on Mathematics Teaching and Learning}, edited by Frank K. Lester, pp. 669--705. Charlotte, N.C.: Information Age Publishing; Reston, Va.: National Council of Teachers of Mathematics, 2007.\vspace{.25\baselineskip}

          \par\noindent\hspace*{1em} This review of literature on early algebraic learning and reasoning discusses the important relationship between arithmetic and elementary algebra. Findings in the reviewed research include the implication that arithmetic can be viewed as a part of algebra as opposed to a separate entity and as such is promising as an entry point into learning algebra. In particular, understanding the structure of numbers, such as numerical properties, may prevent some common misunderstandings in algebra on the part of students.\vspace{.35\baselineskip}

        \item Lehrer, Richard. ``Developing Understanding of Measurement.'' In \emph{A Research Companion to ``Principles and Standards for School Mathematics''}, edited by Jeremy Kilpatrick, W. Gary Martin, and Deborah Schifter, pp. 179--92. Reston, Va.: National Council of Teachers of Mathematics, 2003.\vspace{.25\baselineskip}

          \par\noindent\hspace*{1em} In this survey of research on measurement, the author discusses the centrality of error to measure and approximation. Although much of the research described deals with younger children, it supports the importance of giving students opportunities in the classroom to truggle with and understand error.\vspace{.35\baselineskip}

        \item Sowder, Judith. ``Estimation and Number Sense.'' In \emph{Handbook of Research on Mathematics Teaching and Learning}, edited by Douglas A. Grouws, pp. 371--89. Reston, Va.: National Council of Teachers of Mathematics, 1992.\vspace{.25\baselineskip}

          \par\noindent\hspace*{1em} This review of research on mathematical estimation and number sense discusses three related areas: computational estimation, measurement estimation, and number sense. The chapter makes a case for the importance of both school-aged children's and adults' being able to estimate in both in-school and out-of-school contexts. Posessing a deep understanding of number, including understanding magnitude and numerical comparisons, and being about to reason about number are vital to being a good estimator.\vspace{.35\baselineskip}

        \item National Mathematics Advisory Panel. \emph{Foundations for Success: The Final Report of the National Mathematics Advisory Panel}. http://www.ed.gov/about/bdscomm/list/mathpanel/index.html.\vspace{.25\baselineskip}

          \par\noindent\hspace*{1em} The National Mathematics Advisory Panel was formed by presidential order to explore available research on how to improve mathematics performance among American students, particularly focusing on algebra as a gateway to success in mathematics. Among other recommendations in its final report, the panel highlights the importance of number sense for older students as well as younger students. Number-sense concepts, including place value and magnitude of numbers, as well as how these concepts extend beyond whole numbers to numbers expressed as fractions, decimals, and exponents, are crucial for problem solving. In addition, the report emphasizes the importance of fractions in students' success in algebra.
      \end{enumerate}}
  }
  {
    \VI{
      \begin{enumerate}
        \item Carraher, David W., and Analucia D. Schliemann, ``Early Algebra and Algebraic Reasoning.'' In \emph{Second Handbook of Research on Mathematics Teaching and Learning}, edited by Frank K. Lester, pp. 669--705. Charlotte, N.C.: Information Age Publishing; Reston, Va.: National Council of Teachers of Mathematics, 2007.\vspace{.25\baselineskip}

          \par\noindent\hspace*{1em} Bài tổng quan tài liệu nghiên cứu về việc học và luận đại số ở giai đoạn sớm này thảo luận mối quan hệ quan trọng giữa số học và đại số sơ cấp. Các phát hiện trong nghiên cứu được tổng quan bao gồm hàm ý rằng số học có thể được xem như một phần của đại số, thay vì như thực thể tách biệt; với tư cách đó, số học có triển vọng trở thành điểm vào cho việc học đại số. Đặc biệt, việc hiểu cấu trúc của số, chẳng hạn tính chất số, có thể giúp ngăn ngừa một số hiểu lầm thường gặp của học sinh trong đại số.\vspace{.35\baselineskip}

        \item Lehrer, Richard. ``Developing Understanding of Measurement.'' In \emph{A Research Companion to ``Principles and Standards for School Mathematics''}, edited by Jeremy Kilpatrick, W. Gary Martin, and Deborah Schifter, pp. 179--92. Reston, Va.: National Council of Teachers of Mathematics, 2003.\vspace{.25\baselineskip}

          \par\noindent\hspace*{1em} Trong bài khảo sát nghiên cứu về đo lường này, tác giả thảo luận vai trò trung tâm của sai số trong đo lường và xấp xỉ. Mặc dù phần lớn nghiên cứu được mô tả liên quan đến trẻ nhỏ, những nghiên cứu ấy vẫn ủng hộ tầm quan trọng của việc tạo cơ hội trong lớp học để học sinh vật lộn với sai số và hiểu sai số.\vspace{.35\baselineskip}

        \item Sowder, Judith. ``Estimation and Number Sense.'' In \emph{Handbook of Research on Mathematics Teaching and Learning}, edited by Douglas A. Grouws, pp. 371--89. Reston, Va.: National Council of Teachers of Mathematics, 1992.\vspace{.25\baselineskip}

          \par\noindent\hspace*{1em} Bài tổng quan nghiên cứu về ước lượng toán học và cảm thức số này thảo luận ba mảng liên quan: ước lượng tính toán, ước lượng đo lường và cảm thức số. Chương này lập luận về tầm quan trọng của khả năng ước lượng ở cả trẻ em trong độ tuổi đi học lẫn người lớn, cả trong và ngoài nhà trường. Hiểu sâu về số, bao gồm hiểu độ lớn và so sánh số, cùng với khả năng luận về số, là điều thiết yếu để trở thành người ước lượng tốt.\vspace{3.35\baselineskip}

        \item National Mathematics Advisory Panel. \emph{Foundations for Success: The Final Report of the National Mathematics Advisory Panel}. http://www.ed.gov/about/bdscomm/list/mathpanel/index.html.\vspace{.25\baselineskip}

          \par\noindent\hspace*{1em} Hội đồng Cố vấn Toán học Quốc gia được thành lập theo sắc lệnh của tổng thống nhằm khảo sát nghiên cứu hiện có về cách cải thiện thành tích toán của học sinh tại Hoa Kì, đặc biệt tập trung vào đại số như cửa ngõ dẫn đến thành công trong toán. Bên cạnh các khuyến nghị khác trong báo cáo cuối cùng, hội đồng nhấn mạnh tầm quan trọng của cảm thức số đối với học sinh các lớp trên cũng như học sinh các lớp dưới. Các khái niệm thuộc cảm thức số, bao gồm giá trị theo vị trí và độ lớn của số, cũng như cách các khái niệm này mở rộng vượt khỏi số nguyên không âm sang số được biểu diễn dưới dạng phân số, số thập phân và lũy thừa, là then chốt đối với giải quyết vấn đề. Ngoài ra, báo cáo nhấn mạnh tầm quan trọng của phân số đối với thành công của học sinh trong đại số.
      \end{enumerate}}
  }{\NOTE{}}

\ParallelSection
  {Chapter 5: Reasoning with Algebraic Symbols}
  {Chương 5: Luận với kí hiệu đại số}
  {}

\TriPara
  {
    \EN{
      \begin{itemize}
        \item Fey, James T., ed. \emph{Computing and Mathematics: The Impact on the Secondary School Curricula}. College Park, Md.: University of Maryland, 1984.\vspace{.25\baselineskip}

          \par\noindent\hspace*{1em} This examination of technology in mathematics education from the 1980s includes considerations and comparisons that are still important today. In addition to examining the impact of modern mathematical notation, the authors note that mathematics education influences the discovery of new technology, which in turn influences mathematics education. They describe the history of mathematics education as characterized by cycles of stability and unrest. The editor also notes aspects of the potential impact of technology on curriculum, teaching, and learning.\vspace{.35\baselineskip}

        \item Radford, Luis, and Luis Puig. ``Syntax and Meaning as Sensuous, Visual, Historical Forms of Algebraic Thinking.'' \emph{Educational Studies in Mathematics} 66, no. 2 (October 2007): 145--64.\vspace{.25\baselineskip}

          \par\noindent\hspace*{1em} This article notes that the development of symbolic algebra is considered a great cultural accomplishment. It examines the development of algebraic symbolism from a semiotic viewpoint and observes the effects of culture on both the development of ninth-century Arabic notation and contemporary mathematics students in a modern school setting. The authors highlight some conceptual challenges that can arise in the learning of algebra. One such point is that in learning to understand algebra, students must be able to make some sense of the visual images created by algebraic symbols; doing so may involve understanding multilayered meanings.\vspace{.35\baselineskip}

        \item Saul, Mark. ``Algebra: What Are We Teaching?'' In \emph{The Roles of Representation in School Mathematic}s, 2001 Yearbook of the National Council of Teachers of Mathematics (NCTM), edited by Albert A. Cuoco, pp. 35--43. Reston, Va.: NCTM, 2001.\vspace{.25\baselineskip}

          \par\noindent\hspace*{1em} Saul identifies three levels in learning algebra: generalization of arithmetic, attention to binary operations, and recognition of algebraic form (in which students can conceive of variables as representing complex algebraic expressions). Students with insufficient understanding will cling to memorized rules. Taking away the possibility of trial and error can force students to focus on the operations involved in an algebraic problem and move them from the first level of understanding to the second.\vspace{.35\baselineskip}
        
        \item Kaput, James J., Maria L. Blanton, and Luis Moreno. ``Algebra from a Symbolization Point of View.'' In \emph{Algebra in the Early Grades}, edited by James J. Kaput, Daniel W. Carraher, and Maria L. Blanton, pp. 19--51. New York: Lawrence Erlbaum Associates, 2008.\vspace{.25\baselineskip}

          \par\noindent\hspace*{1em} This analysis of the process of symbolization as it relates to the development of algebraic understanding notes that symbolization is driven by representational economy and communicative power. Communicative power can be measured in terms of argument and justification. When these processes are not present and ``the symbolization process is cut short,'' algebraic learning difficulties can be expected to arise (p. 46). The topics discussed include the transition from arithmetic to algebra, joint variations and functions, and modeling.\vspace{.35\baselineskip}

        \item Katz, Victor J., and Bill Barton. ``Stages in the History of Algebra with Implications for Teaching.'' \emph{Educational Studies in Mathematics} 66, no. 2 (October 2007): 185--201.\vspace{.25\baselineskip}

          \par\noindent\hspace*{1em} The authors include the geometric stage as one of the conceptual stages of the development of algebraic ideas. They note that an introduction to the idea of function through a geometric tool could be a useful approach and that the idea of function was originally developed through geometry. Other ideas include the remaining three stages of conceptual development: static equation-solving, dynamic function, and abstraction. The authors suggest that because algebra arose from the need to solve problems, a problem-solving approach may help overcome the barrier that algebra often presents.
      \end{itemize}}
  }
  {
    \VI{
      \begin{itemize}
        \item Fey, James T., ed. \emph{Computing and Mathematics: The Impact on the Secondary School Curricula}. College Park, Md.: University of Maryland, 1984.\vspace{.25\baselineskip}

          \par\noindent\hspace*{1em} Công trình khảo sát công nghệ trong giáo dục toán từ thập niên 1980 này bao gồm những cân nhắc và so sánh vẫn còn quan trọng đến nay. Bên cạnh việc xem xét tác động của kí pháp toán học hiện đại, các tác giả lưu ý rằng giáo dục toán ảnh hưởng đến việc phát kiến công nghệ mới, và đến lượt mình, công nghệ lại ảnh hưởng đến giáo dục toán. Họ mô tả lịch sử giáo dục toán như được đặc trưng bởi các chu kì ổn định và bất ổn. Biên tập viên cũng lưu ý những khía cạnh trong tác động tiềm tàng của công nghệ đối với chương trình học, dạy học và học tập.\vspace{.35\baselineskip}

        \item Radford, Luis, and Luis Puig. ``Syntax and Meaning as Sensuous, Visual, Historical Forms of Algebraic Thinking.'' \emph{Educational Studies in Mathematics} 66, no. 2 (October 2007): 145--64.\vspace{.25\baselineskip}

          \par\noindent\hspace*{1em} Bài viết này lưu ý rằng sự phát triển của đại số kí hiệu được xem là thành tựu văn hóa lớn. Bài viết khảo sát sự phát triển của hệ kí hiệu đại số từ góc nhìn kí hiệu học, đồng thời quan sát tác động của văn hóa đối với cả sự phát triển của kí pháp Ả Rập thế kỉ thứ chín lẫn học sinh toán đương đại trong môi trường trường học hiện đại. Các tác giả làm nổi bật một số thách thức khái niệm có thể nảy sinh trong việc học đại số. Một điểm như vậy là: khi học để hiểu đại số, học sinh phải có khả năng làm cho những hình ảnh trực quan do kí hiệu đại số tạo ra trở nên có nghĩa ở mức nào đó; điều đó có thể đòi hỏi hiểu những tầng nghĩa chồng lớp.\vspace{1.35\baselineskip}

        \item Saul, Mark. ``Algebra: What Are We Teaching?'' In \emph{The Roles of Representation in School Mathematic}s, 2001 Yearbook of the National Council of Teachers of Mathematics (NCTM), edited by Albert A. Cuoco, pp. 35--43. Reston, Va.: NCTM, 2001.\vspace{.25\baselineskip}

          \par\noindent\hspace*{1em} Saul xác định ba cấp độ trong việc học đại số: khái quát hóa số học, chú ý đến phép toán hai ngôi, và nhận diện dạng đại số, trong đó học sinh có thể quan niệm biến là đại diện cho các biểu thức đại số phức tạp. Học sinh có sự hiểu chưa đầy đủ sẽ bám vào quy tắc ghi nhớ. Việc loại bỏ khả năng thử và sai có thể buộc học sinh tập trung vào các phép toán liên quan trong bài toán đại số và đưa các em từ cấp độ thứ nhất của sự hiểu sang cấp độ thứ hai.\vspace{1.35\baselineskip}
        
        \item Kaput, James J., Maria L. Blanton, and Luis Moreno. ``Algebra from a Symbolization Point of View.'' In \emph{Algebra in the Early Grades}, edited by James J. Kaput, Daniel W. Carraher, and Maria L. Blanton, pp. 19--51. New York: Lawrence Erlbaum Associates, 2008.\vspace{.25\baselineskip}

          \par\noindent\hspace*{1em} Phân tích này về quá trình kí hiệu hóa trong quan hệ với sự phát triển hiểu biết đại số lưu ý rằng kí hiệu hóa được thúc đẩy bởi tính tiết kiệm biểu diễn và sức mạnh giao tiếp. Sức mạnh giao tiếp có thể được đánh giá qua biện giải và biện minh. Khi những quá trình này không hiện diện và ``quá trình kí hiệu hóa bị cắt ngắn'', có thể dự đoán khó khăn trong học đại số sẽ nảy sinh (trang 46). Chủ đề được bàn đến bao gồm chuyển tiếp từ số học sang đại số, biến thiên đồng thời và hàm số, cùng mô hình hóa.\vspace*{1.35\baselineskip}

        \item Katz, Victor J., and Bill Barton. ``Stages in the History of Algebra with Implications for Teaching.'' \emph{Educational Studies in Mathematics} 66, no. 2 (October 2007): 185--201.\vspace{.25\baselineskip}

          \par\noindent\hspace*{1em} Các tác giả xem giai đoạn hình học là một trong các giai đoạn khái niệm trong sự phát triển ý tưởng đại số. Họ lưu ý rằng việc giới thiệu ý tưởng hàm số bằng công cụ hình học có thể là cách tiếp cận hữu ích, và ý tưởng hàm số ban đầu được phát triển thông qua hình học. Các ý tưởng khác được bàn đến bao gồm ba giai đoạn còn lại của phát triển khái niệm: giải phương trình tĩnh, hàm số động và trừu tượng hóa. Các tác giả cho rằng vì đại số nảy sinh từ nhu cầu giải quyết vấn đề, cách tiếp cận giải quyết vấn đề có thể giúp vượt qua rào cản mà đại số thường đặt ra.
      \end{itemize}}
  }{\NOTE{}}

\ParallelSection
  {Chapter 6: Reasoning with Functions}
  {Chương 6: Luận với hàm số}
  {}

\TriPara
  {
    \EN{
      \begin{enumerate}
        \item Yerushalmy, Michal, and Beba Shternberg. ``Charting a Visual Course to the Concept of Function.'' In \emph{The Roles of Representation in School Mathematics}, 2001 Yearbook of the National Council of Teachers of Mathematics (NCTM), edited by Albert A. Cuoco, pp. 251--67. Reston, Va.: NCTM, 2001.\vspace{.25\baselineskip}

          \par\noindent\hspace*{1em} In making the case for the development and use of technology that helps students build a visual understanding of functions, the researchers show that students with a good concept of function are able to solve real-world problems more easily and identify important properties of mathematical representations of situations. They maintain that introduction to symbolic representations can follow the use of other visual representations of functions. The software that they developed, Function Sketcher, allows students to draw the movements in a real-life situation with freehand mouse movements connected to a Cartesian graph, use iconic graph portions for different types of rates of change (increasing, decreasing, constant), and represent incremental rates of change by using stairstep visuals with graphs.\vspace{.35\baselineskip}

        \item Lloyd, Gwendolyn M., and Marvin (Skip) Wilson. ``Supporting Innovation: The Impact of a Teacher's Conceptions of Function on His Implementation of a Reform Curriculum.'' \emph{Journal for Research in Mathematics Education} 29, no. 3 (May 1998): 248--74.\vspace{.25\baselineskip}

          \par\noindent\hspace*{1em} This examination of the effect of one teacher's conception of function on his teaching contains an introductory discussion of the concept of function that notes the powerful and permeating nature of function in mathematics and its ability to give meaning to complex situations. Among the topics discussed are concept images and the power that comes from understanding multiple representations. Each representational format has different strengths in different situations, so the user needs to have an integrated concept image that allows the knowledgeable choice of which one to use. The teacher's dialogues with students reflect an integrated understanding of the function concept that allows him to use reform-based published materials effectively.\vspace{.35\baselineskip}

        \item Coulombe, Wendy N., and Sarah B. Berenson. ``Representations of Patterns and Functions: Tools for Learning.'' In \emph{The Roles of Representation in School Mathematics}, 2001 Yearbook of the National Council of Teachers of Mathematics (NCTM), edited by Albert A. Cuoco, pp. 166--72. Reston, Va.: NCTM, 2001.\vspace{.25\baselineskip}

          \par\noindent\hspace*{1em} This article contains three examples of functions that address different forms of mathematical knowledge. A weight-loss example addresses graphical interpretation and data generation. An iced-tea example addresses verbal pattern description. An allowance example addresses qualitative graphical construction. The authors emphasize the point that the ability to interpret and translate representations can help students construct mental images of patterns and functions and thus extend their algebraic thinking. A deeper interpretation based on familiar events and problem solving can broaden students' understanding of conventional representations beyond mere manipulation. Furthermore, asking a student to interpret a concept by using a representation different from the one with which it is initially presented is one way to determine students' understanding.\vspace{.35\baselineskip}

        \item Chazan, Daniel, and Michal Yerushalmy. “On Appreciating the Cognitive Complexity of School Algebra: Research on Algebra Learning and Directions of Curricular Change.” In \emph{A Research Companion to ``Principles and Standards for School Mathematics''}, edited by Jeremy Kilpatrick, W. Gary Martin, and Deborah Schifter, pp. 123--33. Reston, Va.: National Council of Teachers of Mathematics, 2003.\vspace{.25\baselineskip}

          \par\noindent\hspace*{1em} The authors note that learners experience complexities of thought that the traditional focus on solution methods does not directly address. For example, more differences exist among strings of symbols labeled as equations than the standard definition of equations indicates. More specific terminology might label some equations as formulas, open sentences, identities, functions, or properties, depending on the different uses of variables. To address the need that students have for greater understanding of algebraic symbolization, they should work on tasks that can be approached from multiple perspectives. Content should focus on a small set of big ideas, and instruction should connect with students' experiences.
      \end{enumerate}
    }
  }
  {
    \VI{
      \begin{enumerate}
        \item Yerushalmy, Michal, and Beba Shternberg. ``Charting a Visual Course to the Concept of Function.'' In \emph{The Roles of Representation in School Mathematics}, 2001 Yearbook of the National Council of Teachers of Mathematics (NCTM), edited by Albert A. Cuoco, pp. 251--67. Reston, Va.: NCTM, 2001.\vspace{.25\baselineskip}

          \par\noindent\hspace*{1em} Khi biện minh cho việc phát triển và sử dụng công nghệ giúp học sinh xây dựng hiểu biết trực quan về hàm số, các nhà nghiên cứu chỉ ra rằng học sinh có khái niệm hàm số tốt có thể giải quyết vấn đề thực tế dễ dàng hơn và nhận diện được tính chất quan trọng trong biểu diễn toán học của tình huống. Họ cho rằng việc giới thiệu biểu diễn kí hiệu có thể theo sau việc sử dụng biểu diễn trực quan khác của hàm số. Phần mềm do họ phát triển, Function Sketcher, cho phép học sinh vẽ chuyển động trong tình huống đời thực bằng nét chuột tự do liên kết với đồ thị trong hệ tọa độ Descartes, sử dụng phần đồ thị dạng biểu tượng cho các kiểu tốc độ biến thiên khác nhau (tăng, giảm, không đổi), và biểu diễn tốc độ biến thiên theo từng bước tăng bằng hình bậc thang trên đồ thị.\vspace{1.35\baselineskip}

        \item Lloyd, Gwendolyn M., and Marvin (Skip) Wilson. ``Supporting Innovation: The Impact of a Teacher's Conceptions of Function on His Implementation of a Reform Curriculum.'' \emph{Journal for Research in Mathematics Education} 29, no. 3 (May 1998): 248--74.\vspace{.25\baselineskip}

          \par\noindent\hspace*{1em} Nghiên cứu này xem xét tác động của quan niệm về hàm số ở một giáo viên đối với việc dạy của người đó, đồng thời có phần thảo luận mở đầu về khái niệm hàm số, trong đó lưu ý tính chất mạnh mẽ và lan tỏa của hàm số trong toán học cũng như khả năng đem lại ý nghĩa cho những tình huống phức tạp. Các chủ đề được bàn tới bao gồm hình ảnh khái niệm và sức mạnh đến từ việc hiểu nhiều biểu diễn. Mỗi dạng biểu diễn có thế mạnh khác nhau trong những tình huống khác nhau, vì vậy người sử dụng cần có hình ảnh khái niệm tích hợp, cho phép lựa chọn có hiểu biết về dạng biểu diễn cần dùng. Đối thoại giữa giáo viên và học sinh phản ánh hiểu biết tích hợp về khái niệm hàm số, cho phép giáo viên sử dụng hiệu quả các tài liệu đã xuất bản theo định hướng cải cách.\vspace{.35\baselineskip}

        \item Coulombe, Wendy N., and Sarah B. Berenson. ``Representations of Patterns and Functions: Tools for Learning.'' In \emph{The Roles of Representation in School Mathematics}, 2001 Yearbook of the National Council of Teachers of Mathematics (NCTM), edited by Albert A. Cuoco, pp. 166--72. Reston, Va.: NCTM, 2001.\vspace{.25\baselineskip}

          \par\noindent\hspace*{1em} Bài báo này trình bày ba ví dụ về hàm số, xử lí những dạng tri thức toán học khác nhau. Ví dụ giảm cân xử lí diễn giải bằng đồ thị và tạo dữ liệu. Ví dụ trà đá xử lí mô tả mẫu hình bằng lời. Ví dụ tiền tiêu vặt xử lí xây dựng đồ thị định tính. Các tác giả nhấn mạnh rằng khả năng diễn giải và chuyển đổi biểu diễn có thể giúp học sinh xây dựng hình ảnh tinh thần về mẫu hình và hàm số, từ đó mở rộng tư duy đại số. Diễn giải sâu hơn dựa trên sự kiện quen thuộc và giải quyết vấn đề có thể mở rộng hiểu biết của học sinh về biểu diễn quy ước vượt khỏi thao tác đơn thuần. Ngoài ra, yêu cầu học sinh diễn giải khái niệm bằng biểu diễn khác với biểu diễn ban đầu là một cách xác định sự hiểu của học sinh.\vspace{5.35\baselineskip}

        \item Chazan, Daniel, and Michal Yerushalmy. “On Appreciating the Cognitive Complexity of School Algebra: Research on Algebra Learning and Directions of Curricular Change.” In \emph{A Research Companion to ``Principles and Standards for School Mathematics''}, edited by Jeremy Kilpatrick, W. Gary Martin, and Deborah Schifter, pp. 123--33. Reston, Va.: National Council of Teachers of Mathematics, 2003.\vspace{.25\baselineskip}

          \par\noindent\hspace*{1em} Các tác giả lưu ý rằng người học gặp những phức tạp về tư duy mà trọng tâm truyền thống vào phương pháp giải không trực tiếp xử lí. Chẳng hạn, giữa các chuỗi kí hiệu được gắn nhãn là phương trình có nhiều khác biệt hơn so với điều định nghĩa chuẩn về phương trình thể hiện. Tùy cách dùng biến khác nhau, thuật ngữ cụ thể hơn có thể gọi một số phương trình là công thức, mệnh đề mở, đồng nhất thức, hàm số hoặc tính chất. Để đáp ứng nhu cầu của học sinh về sự hiểu sâu hơn đối với kí hiệu hóa đại số, các em cần làm việc với nhiệm vụ có thể được tiếp cận từ nhiều góc nhìn. Nội dung nên tập trung vào một số ít ý tưởng lớn, và việc dạy cần kết nối với trải nghiệm của học sinh.
      \end{enumerate}}
  }{\NOTE{}}

\ParallelSection
  {Chapter 7: Reasoning with Geometry}
  {Chương 7: Luận với hình học}
  {}

\TriPara
  {
    \EN{
      \begin{enumerate}
        \item Battista, Michael T. ``The Development of Geometric and Spatial Thinking.'' In \emph{Second Handbook of Research on Mathematics Teaching and Learning}, edited by Frank K. Lester, pp. 843--908. Charlotte, N.C.: Information Age Publishing; Reston, Va.: National Council of Teachers of Mathematics, 2007.\vspace{.25\baselineskip}

          \par\noindent\hspace*{1em} In this review of research, the author looks broadly at research on geometric and spatial reasoning. He examines several theories related to students' thinking in geometry and discusses in detail the van Hiele levels, a theory of the way students progress through geometric thinking. The author considers the compelling research supporting these levels of development as informative for understanding students' thinking in geometry.\vspace{1.35\baselineskip}

        \item Clements, Douglas. ``Teaching and Learning Geometry.'' In \emph{A Research Companion to ``Principles and Standards for School Mathematics''}, edited by Jeremy Kilpatrick, W. Gary Martin, and Deborah Schifter, pp. 151--78. Reston, Va.: National Council of Teachers of Mathematics, 2003.\vspace{.25\baselineskip}

          \par\noindent\hspace*{1em} The author discusses research on teaching and learning geometry. In the United States, as compared with other countries, students have fewer meaningful opportunities to study geometry throughout their K-12 education. Providing opportunities for reasoning with geometry at earlier grades would allow high school students' experiences with geometry to be more profound. A discussion of theories of geometrical thinking, research on using technology in geometry, and implications for teaching and curricular practices are included.\vspace{.35\baselineskip}
        
        \item Herbst, Patricio G. ``Engaging Students in Proving: A Double Bind on the Teacher.'' \emph{Journal for Research in Mathematics Education} 33 (May 2002): 176--203.\vspace{.25\baselineskip}

          \par\noindent\hspace*{1em} The author uses data from a high school mathematics classroom to discuss the impact of the traditional, formal two-column proof on the teacher. In particular, he describes a struggle between competing goals for students-developing the ideas for the proof versus proving a proposition. He discusses implications for instruction and the role of formal proof in high school mathematics.
      \end{enumerate}}
  }
  {
    \VI{
      \begin{enumerate}
        \item Battista, Michael T. ``The Development of Geometric and Spatial Thinking.'' In \emph{Second Handbook of Research on Mathematics Teaching and Learning}, edited by Frank K. Lester, pp. 843--908. Charlotte, N.C.: Information Age Publishing; Reston, Va.: National Council of Teachers of Mathematics, 2007.\vspace{.25\baselineskip}

          \par\noindent\hspace*{1em} Trong bài tổng quan nghiên cứu này, tác giả khảo sát rộng các công trình về luận hình học và luận không gian. Ông xem xét nhiều lí thuyết liên quan đến tư duy của học sinh trong hình học và thảo luận chi tiết các cấp độ van Hiele, một lí thuyết mô tả cách học sinh tiến triển qua tư duy hình học. Tác giả xem những nghiên cứu thuyết phục ủng hộ các cấp độ phát triển này là nguồn thông tin hữu ích để hiểu tư duy của học sinh trong hình học.\vspace{.35\baselineskip}

        \item Clements, Douglas. ``Teaching and Learning Geometry.'' In \emph{A Research Companion to ``Principles and Standards for School Mathematics''}, edited by Jeremy Kilpatrick, W. Gary Martin, and Deborah Schifter, pp. 151--78. Reston, Va.: National Council of Teachers of Mathematics, 2003.\vspace{.25\baselineskip}

          \par\noindent\hspace*{1em} Tác giả thảo luận các nghiên cứu về dạy và học hình học. Tại Hoa Kì, so với nhiều quốc gia khác, học sinh có ít cơ hội học hình học một cách có ý nghĩa trong suốt giáo dục K--12. Việc tạo cơ hội cho học sinh luận với hình học ở các lớp sớm hơn sẽ giúp trải nghiệm hình học của học sinh trung học phổ thông trở nên sâu sắc hơn. Chương này cũng bao gồm thảo luận về các lí thuyết tư duy hình học, nghiên cứu về sử dụng công nghệ trong hình học, cũng như hàm ý đối với giảng dạy và thực hành chương trình học.\vspace{1.35\baselineskip}
        
        \item Herbst, Patricio G. ``Engaging Students in Proving: A Double Bind on the Teacher.'' \emph{Journal for Research in Mathematics Education} 33 (May 2002): 176--203.\vspace{.25\baselineskip}

          \par\noindent\hspace*{1em} Tác giả sử dụng dữ liệu từ một lớp toán trung học phổ thông để thảo luận tác động của hình thức chứng minh hai cột truyền thống, mang tính hình thức, đối với giáo viên. Cụ thể, ông mô tả sự giằng co giữa những mục tiêu cạnh tranh dành cho học sinh: phát triển ý tưởng cho chứng minh và chứng minh một mệnh đề. Tác giả cũng bàn về hàm ý đối với dạy học và vai trò của chứng minh hình thức trong toán trung học phổ thông.
      \end{enumerate}}
  }{\NOTE{}}

\ParallelSection
  {Chapter 8: Reasoning with Statistics and Probability}
  {Chương 8: Luận với thống kê và xác suất}
  {}

\TriPara
  {
    \EN{
      \begin{enumerate}
        \item Ganter, Susan L., and William Barker, eds. \emph{The Curriculum Foundations Project: Voices of the Partner Disciplines}. Washington, D.C.: Mathematical Association of America, 2004.\vspace{.25\baselineskip}

          \par\noindent\hspace*{1em} This report describes the work of the Mathematical Association of America's Curriculum Foundations Project, in which representatives from college disciplines outside of mathematics met to discuss their students' mathematical needs with representatives from mathematics departments. The result is a common vision for the first two years of undergraduate mathematics. Statistics and data analysis are noted as crucial for students majoring in many of these disciplines and as such are highlighted as an important part of the undergraduate mathematics experience.\vspace{.35\baselineskip}
        
        \item Franklin, Christine, Gary Kader, Denise Mewborn, Jerry Moreno, Roxy Peck, Mike Perry, and Richard Scheaffer. \emph{Guidelines for Assessment and Instruction in Statistics Education (GAISE) Report: A Pre-K-12 Curriculum Framework}. Alexandria, Va.: American Statistical Association, 2007.\vspace{.25\baselineskip}

          \par\noindent\hspace*{1em} In this framework for the teaching of statistics, leaders in statistics and statistics education make the case for a stronger emphasis on statistics in the K-12 classroom on the basis of the increased necessity for understanding of and the ability to use data. Throughout their school experience, all students should have experiences with, and become proficient in, statistical problem solving. This document fosters greater insight into what these experiences should entail.\vspace{.35\baselineskip}
        
        \item  Kader, Gary, and Jim Mamer. ``Statistics in the Middle Grades: Understanding Center and Spread.'' \emph{Mathematics Teaching in the Middle School} 14 (August 2008): 38--43.\vspace{.25\baselineskip}

          \par\noindent\hspace*{1em} This article, one in a series of articles published for each grade level, is based on the GAISE report of the American Statistical Association (Franklin et al. 2007) and describes the development of basic understanding of the distribution of a data-based variable from elementary to high school. In this article about middle school, the focus is on understanding center and spread. The authors discuss several different mathematical tasks used in classrooms, as well as students' responses.\vspace{.35\baselineskip}
        
        \item Scheaffer, Richard, and Josh Tabor. ``Statistics in the High School Mathematics Curriculum: Building Sound Reasoning under Uncertain Conditions.'' \emph{Mathematics Teacher} 102 (August 2008): 56--61.\vspace{.25\baselineskip}

          \par\noindent\hspace*{1em} This article, one in a series of articles published for each grade level, is based on the GAISE report of the American Statistical Association (Franklin et al. 2007) and describes the development of basic understanding of the distribution of a data-based variable from elementary to high school. In this article about high school, the focus is on building sound reasoning under conditions of uncertainty.
      \end{enumerate}}
  }
  {
    \VI{
      \begin{enumerate}
        \item Ganter, Susan L., and William Barker, eds. \emph{The Curriculum Foundations Project: Voices of the Partner Disciplines}. Washington, D.C.: Mathematical Association of America, 2004.\vspace{.25\baselineskip}

          \par\noindent\hspace*{1em} Báo cáo này mô tả công việc của Dự án Nền tảng Chương trình học thuộc Hiệp hội Toán học Hoa Kì, trong đó đại diện từ các ngành bậc đại học ngoài toán học gặp gỡ đại diện các khoa toán để thảo luận về nhu cầu toán học của sinh viên. Kết quả là hình thành tầm nhìn chung cho hai năm đầu toán bậc đại học. Thống kê và phân tích dữ liệu được ghi nhận là then chốt đối với sinh viên chuyên ngành trong nhiều ngành ấy, và với tư cách đó được nhấn mạnh như phần quan trọng trong trải nghiệm học toán bậc đại học.\vspace{2.35\baselineskip}
        
        \item Franklin, Christine, Gary Kader, Denise Mewborn, Jerry Moreno, Roxy Peck, Mike Perry, and Richard Scheaffer. \emph{Guidelines for Assessment and Instruction in Statistics Education (GAISE) Report: A Pre-K-12 Curriculum Framework}. Alexandria, Va.: American Statistical Association, 2007.\vspace{.25\baselineskip}

          \par\noindent\hspace*{1em} Trong khung dành cho việc dạy thống kê này, các nhà lãnh đạo trong thống kê và giáo dục thống kê đưa ra cơ sở ủng hộ việc nhấn mạnh thống kê nhiều hơn trong lớp học K--12, dựa trên nhu cầu ngày càng tăng về hiểu dữ liệu và khả năng sử dụng dữ liệu. Trong suốt trải nghiệm học tập ở nhà trường, mọi học sinh cần được trải nghiệm và trở nên thành thạo trong giải quyết vấn đề thống kê. Tài liệu này giúp hiểu rõ hơn những trải nghiệm ấy nên bao gồm điều gì.\vspace{.35\baselineskip}
        
        \item  Kader, Gary, and Jim Mamer. ``Statistics in the Middle Grades: Understanding Center and Spread.'' \emph{Mathematics Teaching in the Middle School} 14 (August 2008): 38--43.\vspace{.25\baselineskip}

          \par\noindent\hspace*{1em} Bài báo này, một bài trong chuỗi bài được công bố cho từng cấp lớp, dựa trên báo cáo GAISE của Hiệp hội Thống kê Hoa Kì (Franklin et al. 2007) và mô tả sự phát triển hiểu biết cơ bản về phân bố của biến dựa trên dữ liệu từ tiểu học đến trung học phổ thông. Trong bài viết về trung học cơ sở này, trọng tâm là hiểu trung tâm và độ phân tán. Các tác giả thảo luận một số nhiệm vụ toán học khác nhau được sử dụng trong lớp học, cũng như phản hồi của học sinh.\vspace{1.35\baselineskip}
        
        \item Scheaffer, Richard, and Josh Tabor. ``Statistics in the High School Mathematics Curriculum: Building Sound Reasoning under Uncertain Conditions.'' \emph{Mathematics Teacher} 102 (August 2008): 56--61.\vspace{.25\baselineskip}

          \par\noindent\hspace*{1em} Bài báo này, một bài trong chuỗi bài được công bố cho từng cấp lớp, dựa trên báo cáo GAISE của Hiệp hội Thống kê Hoa Kì (Franklin et al. 2007) và mô tả sự phát triển hiểu biết cơ bản về phân bố của biến dựa trên dữ liệu từ tiểu học đến trung học phổ thông. Trong bài viết về trung học phổ thông này, trọng tâm là xây dựng luận vững chắc trong điều kiện bất định.
      \end{enumerate}}
  }
  {
    \NOTE{Nhóm tài liệu trong Chương 8 cho thấy thống kê và xác suất được đặt ở trung tâm của năng lực học toán trong thế giới dữ liệu. Điểm chung của bốn nguồn không chỉ là bổ sung nội dung thống kê vào chương trình, mà là làm rõ kiểu luận cần hình thành khi học sinh làm việc với dữ liệu, biến thiên, phân bố và bất định.

    \emph{The Curriculum Foundations Project} cho thấy nhu cầu toán học của sinh viên không thể được xác định chỉ từ bên trong khoa toán. Khi đại diện từ ngành đối tác cùng thảo luận với khoa toán, thống kê và phân tích dữ liệu hiện lên như phần then chốt của hai năm đầu toán bậc đại học. Điều này gợi một cách nhìn quan trọng: nền tảng toán học hiện đại không chỉ gồm kĩ thuật hình thức, mà còn gồm năng lực đọc dữ liệu, phân tích biến thiên và lập luận từ bằng chứng định lượng.

    Báo cáo GAISE mở rộng tinh thần ấy xuống giáo dục K--12. Khi nhu cầu hiểu và sử dụng dữ liệu tăng lên, mọi học sinh cần được trải nghiệm giải quyết vấn đề thống kê xuyên suốt quá trình học. Trọng tâm không nằm ở việc học rời rạc vài biểu đồ hay công thức, mà ở việc học sinh biết đặt câu hỏi, thu thập dữ liệu, phân tích, diễn giải và trả lời trên cơ sở dữ liệu.

    Bài của Kader and Mamer làm rõ một mắt xích quan trọng ở trung học cơ sở: hiểu trung tâm và độ phân tán của phân bố dữ liệu. Đây là nền tảng để học sinh không nhìn dữ liệu như danh sách giá trị rời rạc, mà như phân phối có cấu trúc, có xu hướng và có biến thiên. Những nhiệm vụ lớp học và phản hồi của học sinh trong bài giúp thấy luận thống kê được hình thành qua việc so sánh, mô tả và diễn giải dữ liệu.

    Bài của Scheaffer and Tabor đưa trọng tâm lên trung học phổ thông: xây dựng luận vững chắc trong điều kiện bất định. Đây là điểm nối trực tiếp giữa thống kê và xác suất trong FHSM. Học sinh không chỉ tính xác suất hay mô tả dữ liệu, mà cần dùng mô hình xác suất và luận thống kê để đánh giá kết quả trong tình huống có biến thiên và không chắc chắn.

    Nhìn chung, các nguồn này củng cố cách hiểu rằng luận với thống kê và xác suất không phải phần phụ của chương trình toán trung học phổ thông. Nó là nơi học sinh học cách suy nghĩ bằng dữ liệu, nhận ra biến thiên, mô hình hóa bất định và đưa ra kết luận có căn cứ. Theo nghĩa đó, Chương 8 đặt thống kê và xác suất vào đúng vị trí của chúng trong giáo dục toán hiện đại: không chỉ là nội dung cần học, mà là hình thức luận cần thiết để hiểu và hành động trong thế giới dữ liệu.}
  }

\ParallelSection
  {Chapter 9: Equity}
  {Chương 9: Công bằng}
  {}

\TriPara
  {
    \EN{
      \begin{enumerate}
        \item Tate, William, and Celia Rousseau. ``Access and Opportunity: The Political and Social Context of Mathematics Education.'' In \emph{Handbook of International Research in Mathematics Education}, edited by Lyn D. English, pp. 271--99. Mahwah, N.J.: Lawrence Erlbaum Associates, 2002.\vspace{.25\baselineskip}

          \par\noindent\hspace*{1em} In this handbook chapter, the authors look at the political, historical, and social context of mathematics education and the resulting impact on equitable access to school mathematics. Discrepancies in course taking, tracking, teacher quality, and equity at the classroom level are discussed as having an impact on differences in learning opportunities for students of color, students in urban schools, and students of low socioeconomic status.\vspace{1.35\baselineskip}

        \item Tate, William F., and Celia Rousseau. ``Engineering Change in Mathematics Education: Research, Policy, and Practice.'' In \emph{Second Handbook of Research on Mathematics Teaching and Learning}, edited by Frank K. Lester, pp. 1209--46. Charlotte, N.C: Information Age Publishing; Reston, Va.: National Council of Teachers of Mathematics, 2007.\vspace{.25\baselineskip}

          \par\noindent\hspace*{1em} This chapter provides direction and strategies for teachers, schools, and districts trying to improve mathematics teaching and learning in the No Child Left Behind era. As a result, it discusses the current pressures on schools because of the increased focus on testing. In particular, the authors note that the increased focus on testing is detrimental for all students, but in particular for poor and minority students, whose opportunities for important mathematical reasoning and thinking are limited as a result.\vspace{2.35\baselineskip}

        \item Education Trust. \emph{Gaining Traction, Gaining Ground: How Some High Schools Accelerate Learning for Struggling Students}. Washington, D.C.: Education Trust, 2005.\vspace{.25\baselineskip}

          \par\noindent\hspace*{1em} This work is a report of a study comparing seven public high schools-four of which demonstrated incredible growth in students who entered high school significantly behind their peers. These four exemplar schools were compared against high schools that had average growth. The schools all had similar demographics-predominantly minority and low-socioeconomic-status student populations. The study reveals those approaches adopted in the successful schools.\vspace{2.35\baselineskip}

        \item Planty, Michael, Stephen Provasnik, and Bruce Daniel. \emph{High School Coursetaking: Findings from the Condition of Education}, 2007. Washington, D.C.: U.S. Department of Education, National Center for Education Statistics, 2007.\vspace{.25\baselineskip}

          \par\noindent\hspace*{1em} This report provides statistics on trends in United States high school course taking from 1982 through 2005. Some of the trends discussed include the coursework offered by different schools, changes in graduation requirements over the past two decades, and the percent of students who take advanced coursework in science and mathematics. The report also compares trends across racial and socioeconomic subgroups.\vspace{.35\baselineskip}
      \end{enumerate}}
  }
  {
    \VI{
      \begin{enumerate}
        \item Tate, William, and Celia Rousseau. ``Access and Opportunity: The Political and Social Context of Mathematics Education.'' In \emph{Handbook of International Research in Mathematics Education}, edited by Lyn D. English, pp. 271--99. Mahwah, N.J.: Lawrence Erlbaum Associates, 2002.\vspace{.25\baselineskip}

          \par\noindent\hspace*{1em} Trong chương sách này, các tác giả xem xét bối cảnh chính trị, lịch sử và xã hội của giáo dục toán, cũng như tác động của các yếu tố đó đối với việc tiếp cận công bằng môn toán ở trường học. Chênh lệch trong việc lựa chọn môn học, phân luồng, chất lượng giáo viên và mức độ công bằng trong lớp học được thảo luận như yếu tố ảnh hưởng đến sự khác biệt về cơ hội học tập của học sinh thuộc các nhóm thiểu số, học sinh ở khu vực đô thị và học sinh có hoàn cảnh kinh tế - xã hội khó khăn.\vspace{.35\baselineskip}

        \item Tate, William F., and Celia Rousseau. ``Engineering Change in Mathematics Education: Research, Policy, and Practice.'' In \emph{Second Handbook of Research on Mathematics Teaching and Learning}, edited by Frank K. Lester, pp. 1209--46. Charlotte, N.C: Information Age Publishing; Reston, Va.: National Council of Teachers of Mathematics, 2007.\vspace{.25\baselineskip}

          \par\noindent\hspace*{1em} Chương này đưa ra định hướng và chiến lược dành cho giáo viên, nhà trường và học khu đang tìm cách cải thiện việc dạy và học toán trong thời kì No Child Left Behind. Vì vậy, chương bàn đến những áp lực hiện nay mà nhà trường phải đối mặt do việc ngày càng chú trọng kiểm tra đánh giá. Đặc biệt, các tác giả lưu ý rằng sự chú trọng ngày càng tăng vào kiểm tra đánh giá gây bất lợi cho tất cả học sinh, nhưng tác động ấy đặc biệt nặng nề đối với học sinh nghèo và học sinh thuộc nhóm thiểu số, vì nó làm thu hẹp cơ hội để các em tham gia vào hoạt động luận và tư duy toán học quan trọng.\vspace{.35\baselineskip}

        \item Education Trust. \emph{Gaining Traction, Gaining Ground: How Some High Schools Accelerate Learning for Struggling Students}. Washington, D.C.: Education Trust, 2005.\vspace{.25\baselineskip}

          \par\noindent\hspace*{1em} Công trình này là báo cáo nghiên cứu so sánh bảy trường trung học công lập, trong đó bốn trường cho thấy sự tiến bộ vượt bậc của học sinh bước vào trung học với mức độ học tập kém hơn đáng kể so với bạn bè cùng trang lứa. Bốn trường điển hình này được so sánh với các trường có mức tăng trưởng trung bình. Các trường trong nghiên cứu có đặc điểm dân số tương tự nhau, với đa số học sinh thuộc nhóm thiểu số và có hoàn cảnh kinh tế - xã hội khó khăn. Nghiên cứu làm rõ cách tiếp cận đã được áp dụng tại các trường đạt kết quả thành công.\vspace{.35\baselineskip}

        \item Planty, Michael, Stephen Provasnik, and Bruce Daniel. \emph{High School Coursetaking: Findings from the Condition of Education}, 2007. Washington, D.C.: U.S. Department of Education, National Center for Education Statistics, 2007.\vspace{.25\baselineskip}

          \par\noindent\hspace*{1em} Báo cáo này cung cấp số liệu thống kê về xu hướng lựa chọn môn học ở bậc trung học phổ thông tại Hoa Kì trong giai đoạn từ năm 1982 đến năm 2005. Một số xu hướng được thảo luận bao gồm các môn học được cung cấp tại các trường khác nhau, thay đổi trong yêu cầu tốt nghiệp trong hai thập kỉ qua, và tỷ lệ học sinh theo học các môn khoa học và toán học nâng cao. Báo cáo cũng so sánh xu hướng này giữa các nhóm chủng tộc và các nhóm có điều kiện kinh tế - xã hội khác nhau.
      \end{enumerate}}
  }{\NOTE{}}

\ParallelSection
  {Chapter 10: Coherence}
  {Chương 10: Tính mạch lạc}
  {}

\TriPara
  {
    \EN{
      \begin{enumerate}
        \item Black, Paul, and Dylan Wiliam. ``Inside the Black Box: Raising Standards through Classroom Assessment.'' \emph{Phi Delta Kappan} 80 (October 1998): 139--44.\vspace{.25\baselineskip}

          \par\noindent\hspace*{1em} This article looks at research to show that formative assessment has a major impact on classroom instruction. Evidence also shows that formative assessment can have a major impact on students' learning. The authors suggest improvements in formative assessment and ways to make those improvements.\vspace{.35\baselineskip}

        \item National Science Board. \emph{A National Action Plan for Addressing the Critical Needs of the U.S. Science, Technology, Engineering, and Mathematics Education System}. Arlington, Va.: National Science Foundation, 2007.\vspace{.25\baselineskip}

          \par\noindent\hspace*{1em} This report calls attention to the need to address issues in STEM (science, technology, engineering, and mathematics) education. The National Science Board makes strong recommendations for alignment within K-12 STEM education and then between K-12 STEM education and STEM higher education and the workforce. Specific recommendations related to this goal of alignment and others are provided.\vspace{.35\baselineskip}

        \item Smith, Jack, Beth Herbel-Eisenmann, Amanda Jansen, and Jon Star. ``Studying Mathematical Transitions: How Do Students Navigate Fundamental Changes in Curriculum and Pedagogy?'' Paper presented at the 2000 annual meeting of the American Educational Research Association, New Orleans, April 2000.\vspace{.25\baselineskip}

          \par\noindent\hspace*{1em} The authors report on research looking at the impact of transitioning from high school to college and between using traditional curricula and reform curricula. Moving between different curricular approaches is difficult for students because the expectations for learning and doing mathematics are very different with each of these approaches. Students then must then navigate between, in essence, two different mathematical environments.
      \end{enumerate}}
  }
  {
    \VI{
      \begin{enumerate}
        \item Black, Paul, and Dylan Wiliam. ``Inside the Black Box: Raising Standards through Classroom Assessment.'' \emph{Phi Delta Kappan} 80 (October 1998): 139--44.\vspace{.25\baselineskip}

          \par\noindent\hspace*{1em} Bài báo này điểm lại các nghiên cứu để cho thấy đánh giá quá trình có tác động lớn đến hoạt động dạy học trên lớp. Bằng chứng cũng cho thấy đánh giá quá trình có thể tác động mạnh đến việc học của học sinh. Các tác giả đề xuất những hướng cải thiện đánh giá quá trình và cách thực hiện những cải thiện ấy.\vspace{.35\baselineskip}

        \item National Science Board. \emph{A National Action Plan for Addressing the Critical Needs of the U.S. Science, Technology, Engineering, and Mathematics Education System}. Arlington, Va.: National Science Foundation, 2007.\vspace{.25\baselineskip}

          \par\noindent\hspace*{1em} Báo cáo này nhấn mạnh nhu cầu xử lí các vấn đề trong giáo dục STEM (khoa học, công nghệ, kĩ thuật và toán học). Hội đồng Khoa học Quốc gia đưa ra khuyến nghị mạnh mẽ về sự liên kết trong nội bộ giáo dục STEM K--12, rồi giữa giáo dục STEM K--12 với giáo dục đại học STEM và lực lượng lao động STEM. Báo cáo cũng nêu các khuyến nghị cụ thể gắn với mục tiêu liên kết này cùng những mục tiêu khác.\vspace{.35\baselineskip}

        \item Smith, Jack, Beth Herbel-Eisenmann, Amanda Jansen, and Jon Star. ``Studying Mathematical Transitions: How Do Students Navigate Fundamental Changes in Curriculum and Pedagogy?'' Paper presented at the 2000 annual meeting of the American Educational Research Association, New Orleans, April 2000.\vspace{.25\baselineskip}

          \par\noindent\hspace*{1em} Các tác giả báo cáo nghiên cứu xem xét tác động của việc chuyển tiếp từ trung học phổ thông lên đại học, cũng như của việc chuyển đổi giữa chương trình học truyền thống và chương trình học cải cách. Việc di chuyển giữa các cách tiếp cận chương trình học khác nhau gây khó khăn cho học sinh, vì kì vọng về việc học và làm toán trong mỗi cách tiếp cận ấy rất khác nhau. Do đó, về thực chất, học sinh phải điều hướng giữa hai môi trường toán học khác nhau.
      \end{enumerate}}
  }{\NOTE{}}

\CloseGrid